\section{Results}

\begin{figure}[t]
\centering
\includegraphics[width=0.95\linewidth]{fatrop_oval_obstacles_traj.png}
\caption{Representative FATROP solution on the Oval track with obstacles. These feasible, obstacle-avoiding trajectories are the source of the offline expert data used for DT training.}
\label{fig:oval-fatrop}
\end{figure}

\subsection{Offline Training Improved, But Deployment Did Not}

The clearest positive result is that the training pipeline improved substantially over the course of the study. The strongest offline model was the failure-conditioned fine-tune, which reduced validation loss from $0.002835$ for the larger-context parent model to $0.002364$. The lateral-stability fine-tune also improved over the parent to $0.002773$, showing that targeted data and loss design were effective at fitting the expert distribution. Figure~\ref{fig:train-loss} shows a representative training curve for one of the later fine-tuning runs, and Figure~\ref{fig:oval-fatrop} shows the type of feasible obstacle-avoiding trajectories used to generate the offline supervision.

The central limitation is that this offline ranking did not match deployment behavior. The strongest completed warm-start setting before the failure-conditioned run was the post-projection fine-tune with moderate projection, even though it was not the best offline model. This gap between offline imitation quality and warm-start utility is one of the main empirical findings of the report.

\begin{figure}[t]
\centering
\includegraphics[width=0.78\linewidth]{postproj_ft2_train_loss_curves.png}
\caption{Training action-loss curve for the post-projection fine-tune. Training remained stable under the larger-context DT configuration used for the later fine-tuning experiments.}
\label{fig:train-loss}
\end{figure}

\subsection{Best Deployment Setting Still Fell Short of an End-to-End Win}

A completed 10-scenario-per-gate ablation compared the larger-context parent model and the post-projection fine-tune under different levels of post-rollout correction. The best deployment setting among these runs was the post-projection fine-tune with moderate projection. In that setting, the optimizer could sometimes converge slightly faster from the DT proposal than from the baseline initializer. On the 10-scenario obstacle benchmark, for example, baseline solve time was $54.48$s while DT solve time was $52.58$s.

That improvement was not enough to change the end-to-end conclusion. Once warm-start inference time was included, DT total time rose to $55.21$s, and acceptance remained $0\%$ under the default gate. Iteration reduction also remained $0\%$ across the completed default-gate ablation runs. In other words, the learned warm-starts were sometimes promising enough to help the solver locally, but not yet reliable enough to produce a robust full-pipeline gain.

\subsection{Relaxed Acceptance Increased Use of DT Warm-Starts}

After fixing the repeated-start confound, a corrected relaxed-gate diagnostic for the post-projection fine-tune with \texttt{projection=soft} showed that acceptance was not literally impossible. On the corrected 10-scenario evaluation:
\begin{itemize}
\item no-obstacle gate: acceptance $20\%$, baseline solve time $49.06$s, DT total time $51.98$s, iterations $206.8$ vs baseline $191.8$;
\item obstacle gate: acceptance $20\%$, baseline solve time $56.36$s, DT total time $60.79$s, iterations $256.0$ vs baseline $230.7$.
\end{itemize}

These runs show that useful trajectory structure was present in the learned proposals: acceptance rose to $20\%$ on both corrected gates once the rejection logic was relaxed. At the same time, the accepted proposals still did not reduce solver work on average. The implication is that acceptance alone is not the relevant success criterion; accepted warm-starts must also reduce total optimization effort.

\subsection{Final Post-Fix Smoke Benchmark}

The final post-fix smoke rerun focused on the failure-conditioned model, after the key correctness bugs had been fixed. These are small benchmarks with only 3 scenarios per gate and should be interpreted as smoke tests, not definitive leaderboard results. Table~\ref{tab:final-smoke} summarizes them.

\begin{table}[t]
\centering
\caption{Final post-fix smoke benchmark for the failure-conditioned fine-tune with \texttt{projection=soft}. Each gate has only 3 scenarios.}
\label{tab:final-smoke}
\small
\setlength{\tabcolsep}{4pt}
\resizebox{\linewidth}{!}{%
\begin{tabular}{lcccccc}
\toprule
Gate & Baseline Solve & DT Solve & DT Total & Iterations (DT/Base) & WS Accept & Fallback Mean \\
\midrule
No obstacle & 47.91s & 46.34s & 49.04s & 186.3 / 186.3 & 0\% & 0.0 \\
1--4 obstacles & 64.85s & 62.47s & 65.18s & 280.7 / 280.7 & 0\% & 1.0 \\
\bottomrule
\end{tabular}
\par}
\end{table}

These smoke results are consistent with the overall project conclusion. The DT sometimes lowers raw solve time slightly, but once warm-start inference is included, total time is still not better than baseline. More importantly, warm-start acceptance remains $0\%$, and iteration reduction remains $0\%$.

\subsection{Final Smoke Tests Support the Same Picture}

The final failure-conditioned smoke rerun also helps explain why acceptance remained at $0\%$. Projection burden stayed above the default gate: on the no-obstacle gate, mean projection total magnitude was $127.29$ and mean max-step projection was $4.91$; on the obstacle gate, the corresponding values were $121.03$ and $4.90$. This is still a useful result. Earlier failures were often dominated by immediate fallback, whereas the final corrected pipeline produced trajectories that were more structured and more stable, but still not laterally accurate enough to be admitted under the default deployment rule.
